% numbers and result paragraphs; sources: study/insitu/out/server/*.json, docs/CRITIC_INSITU_1.md
\newcommand{\nknotimg}{124}
\newcommand{\splineerr}{0.19}
\newcommand{\gradspeedA}{10}
\newcommand{\iterspeedA}{5}
\newcommand{\breakevenA}{300}
\newcommand{\nseeds}{3}
\newcommand{\gpu}{RTX 3090}
\newcommand{\nbench}{16{,}000}

\newcommand{\speedtext}{%
At that setting the pupil model costs 23\,$\mu$s per emitter and the spline 2.3\,$\mu$s; the fixed cost of the
knot route (one build of the knots that keeps the computation graph, and one backward pass) is 7\,ms. The
gradient is cheaper from about 320 emitters on, 7.6 times at $10^4$ and 10 times at $10^5$ emitters and above.
Re-localizing the emitters costs another 2.5\,$\mu$s per emitter in both routes, so a whole iteration is 5.0 to
5.3 times cheaper. With the coarser 64/3 setting, whose error (0.39\,\%) is twice that of the spline, the
numbers are 4 (gradient) and 2.5 (iteration); a complete fit of 16{,}000 emitters, 40 iterations, took 5.3\,s
through the knots and 8.9\,s per emitter, and both ended at the same Zernike error (Fig.~\ref{fig:speed},
right). The 256/6 setting gives a factor 46 for the gradient at $10^5$ emitters, but that pupil model is ten
times more accurate than the spline it is compared with, so this row is not a like-for-like comparison.}

\newcommand{\gradtext}{%
At the starting point of a fit (Zernikes from the bead stack, or the truth plus 0.01 waves) the two gradients
agree to 0.2\,\% (cosine similarity 1.000000). Near the optimum the gradient is small, and the interpolation
error of the spline becomes visible: at the true parameters the two routes differ by 4\,\% for the Zernikes, 4\,\%
for a pupil phase map and 12\,\% for a pupil amplitude map (measured by the reviewer of this study,
\texttt{critic/grad\_speed.py}, 2000 emitters). On noise-free data this moves the optimum of the Zernike
coefficients by $1.6\times10^{-5}$ waves rms. That is far below any physical relevance, but it is the size of
the statistical error at about $10^6$ emitters, so beyond that a finer spline ($S=3$) or a final step with the
per-emitter route would be needed. In the fits of this study the Zernike error at convergence is set by the
difference between the simulated and the fitted pupil sampling ($5$--$8\times10^{-4}$ waves in scenario A).}

\newcommand{\acctext}{%
Table~\ref{tab:acc} and Fig.~\ref{fig:bias}.

\emph{Beads only.} With the aberration change chosen here, a PSF from the bead stack alone is far off on the
sample data: 15--18 nats per emitter, z errors of 65--80\,nm rms, and for the Zernike model a y error that
changes by 120--180\,nm over the z range. This is a statement about the chosen aberration change, not a general
one, and the bead spline uses the paraxial depth scaling without the empirical correction that is common
practice.

\emph{Scenario A.} The Zernike model refined on the movie patches, starting from the bead fit, is
indistinguishable from the true PSF in every column (8\,s for 11{,}600 patches). Free knots on top add noise
(z rms 22 to 27\,nm, a z-scale error of $-13\pm4$\,nm). A free pupil map on top changes nothing within the
spread over seeds.

\emph{Scenario B.} The refined Zernike model leaves a likelihood deficit of 2.7 nats per emitter, a y error
that grows by $45\pm1$\,nm over the z range, a z-scale error of $19\pm2$\,nm and a photon count that is 15\,\%
low. Free knots on top remove most of the deficit (0.22) and leave the y tilt ($43\pm1$\,nm), the z tilt and the
photon count where they were. This is the prediction of section~\ref{sec:gauge}: the three quantities that
blinking data do not determine stay at the values of the starting model. The free pupil map removes the deficit
(0.05), brings y and z precision to those of the true PSF, and reduces the y tilt to $8\pm5$\,nm on the movie
patches and $18\pm1$\,nm on clean patches; the photon count stays 9--10\,\% low. So the pupil physics removes
most of the z-dependent error but not all of it. Why the residual is larger on clean patches than on movie
patches was not investigated.

\emph{Shared spline.} A single free-knot spline fitted to beads and movie together fits the movie images
(deficit 0.2--0.3) but keeps the geometry of the bead stack: z offset 65\,nm, z-scale error of 75--107\,nm over
the range, photon counts 27--38\,\% low. A model-free spline is therefore not a usable in-situ PSF on its own
when the bead PSF and the sample PSF differ.

The movie gives few emitters at the top of the z range: the detection finds 57--67\,\% of the on-events below
500\,nm, 17\,\% at 600--700\,nm and 4\,\% at 700--800\,nm (Fig.~\ref{fig:zhist}), and the patches come from 700
sites, so they are not 11{,}600 independent emitters.}

\newcommand{\pitfalltext}{%
The spline stores the PSF on a support that allows emitter offsets up to $x_\text{max}$ from the patch centre,
and the localizer clamps the offset at that limit. In the first version of this study $x_\text{max}$ was
1.5 pixels. The patches are centred on the detection peak, and for a defocused, aberrated spot the peak is not
at the emitter: here it moves away by 1.2 pixels on average at 600--700\,nm height and by more than 3 pixels for
some spots. About 3\,\% of the patches then had the emitter beyond the limit, all at large z and on the same
side. The localizer could not reach them, and a flexible PSF model absorbed the difference: the free pupil map
showed a y tilt of 28--38\,nm in scenario A, where the Zernike start was already correct. The 17-term Zernike
model was too stiff to follow and stayed correct, which hid the problem. Controls located it: the same emitters
cropped at their true position gave no tilt, isolated emitters cropped at their detection peak gave 58\,nm, and
with $x_\text{max}=3$ pixels (all results above) the tilt is $1\pm3$\,nm. The independent review found the clamp
as the mechanism. The fraction of fitted emitters at the limit is now reported for every fit (below 0.1\,\%
here). Any localizer with a bounded offset has this failure mode when it is used inside a PSF estimation.}

\newcommand{\patterntext}{%
A lateral shift of the PSF with z is invisible in ordinary blinking data (section~\ref{sec:gauge}). A second,
PSF-independent measurement of the lateral position removes that freedom. We simulated emitters seen in 6
frames under a sinusoidal illumination pattern (2 orientations, 3 phases, pitch 220\,nm, modulation 0.9;
12{,}000 photons in total, 3000 emitters, clean patches); the emitter position is shared by the PSF term and the
pattern term of the likelihood. Table~\ref{tab:pattern}: with the pattern, the free pupil map in scenario B has
mean errors within $\pm1$\,nm in y and $\pm3$\,nm in z over the whole range, against a y error up to 14\,nm from
the same frames without pattern. Free knots lose their y error (44 to 2\,nm) but keep a z error of $-12$ to
$+21$\,nm: a lateral pattern says nothing about the labelling of the z axis.

This is a best case. The pattern (pitch, phases, modulation, registration to the camera) is taken as exactly
known, whereas in an experiment it is estimated from the same localizations; the emitters have three times the
photons of the movie; and the fit has to start from a PSF that is already roughly right (it first refines the
Zernikes on the summed frames), because from the bead PSF emitters lock onto the wrong fringe.}

\newcommand{\limittext}{%
\begin{itemize}\setlength\itemsep{0pt}
\item Simulation only. Data and fit share the vectorial model, NA, refractive indices, wavelength, dipole model
and Poisson noise; only the pupil sampling differs. A check by the reviewer with a sample index of 1.35 instead
of the 1.33 assumed by the fit gave a deficit of 0.10 nats and a z bias of $-22$ to $-5$\,nm for the Zernike
model: the optical constants have to be fitted or known.
\item Three seeds; the aberration change between beads and sample and the content of scenario B are single
arbitrary choices.
\item The speed comparison is between two routes in one code base on one GPU. No comparison with uiPSF or the
code of Droste et al.\ on the same data.
\item Field-constant PSF, uniform background, $15\times15$ patches, isolated spots only.
\item The regularization of the free-knot fits was chosen by looking at the test likelihood of scenario A.
\item The pattern experiment assumes a perfectly known pattern.
\item The settings where the speed matters most (field-dependent or time-dependent PSFs, very large data sets)
are not tested.
\end{itemize}}
